\section{Background}
Student dropout is consequential for both learners and institutions, and its dynamics often become visible through changes in engagement across an academic term. As instruction increasingly shifts to online and hybrid formats, time-stamped learning traces make it possible to examine dropout risk with an explicitly temporal perspective.

\subsection{Why Dropout Matters: Impacts and Scale}
Student dropout has far-reaching consequences at the individual, family, institutional, and societal levels. For individuals, dropout disrupts educational attainment, limits skill development, and can negatively affect motivation and well-being. At the family level, it often represents a substantial sunk financial investment and constrains intergenerational social mobility \cite{Kamissa2020DroppingOut}. Institutions face declining program quality, reputational harm, and revenue losses from reduced enrollment and opportunity costs \cite{Guzman2021DropoutRural, Schneider2010FirstYearAttrition}.

In the United States, approximately 30\% of undergraduates drop out after their first year, resulting in over 9 billion dollars annually in public expenditures with limited social returns \cite{Schneider2010FirstYearAttrition}. Globally, despite expanded access to tertiary education, substantial dropout persists across systems as documented in socio-economic reviews \cite{Aina2022DeterminantsDropout}. Course-level dropout further increases the risk of program withdrawal, institutional exit, and reduced re-enrollment, particularly among marginalized students \cite{McKinney2019CourseDropping, Gicheva2025PersistenceIntroCourse}. These cascading effects underscore the need for early detection and timely support.

\subsection{Structural Shifts in Delivery: From Distance to Online}
Student dropout must be understood within the structural transformation of higher education delivery. Although online learning appears recent, distance education has historical roots in correspondence-based instruction and evolved into online formats with the expansion of the internet \cite{Kentnor2015DistanceEducation}. The COVID-19 pandemic accelerated this shift, forcing institutions worldwide to rapidly adopt online delivery and redefining online education as a core instructional mode rather than a supplement \cite{Bartolic2022COVIDTeaching, Bryson2020RapidAdoption, Zhang2023OnlineLearningReview}.

This transition embedded online education into institutional strategy, led to widespread program migration, and prompted substantial pedagogical redesign. Learning Management Systems (LMS) consequently became central infrastructures for teaching and engagement, generating rich, time-stamped behavioral data \cite{Turnbull2021TransitioningELearning}. These data enable temporal analytics that move dropout modeling beyond static prediction.

Dropout in online higher education reflects a distinct ``dropout ecology'' shaped by social isolation, reduced interaction, and heightened self-regulation demands \cite{Rahmani2024OnlineDropoutReview}. Attrition arises from intersecting academic, institutional, and personal barriers, including workload, financial, health, and technological constraints. While socioeconomic status and academic performance consistently predict early dropout \cite{BarraganMoreno2024ComplexitiesDropout}, comparative analyses across disciplines and study levels remain limited, with much research focused on MOOCs \cite{Prenkaj2021SurveyMLDropout,Sghir2022PredictiveLearningAnalytics,Alalawi2023MLStudentPerformance}.
Learning Management Systems enable detailed tracking of engagement, which correlates with academic outcomes \cite{Mozahem2020LMSActivity}. Although machine-learning models achieve high predictive accuracy, challenges persist due to data heterogeneity, limited reproducibility, and insufficient temporal and contextual modeling \cite{Albreiki2021MLPerformanceReview,Schneider2010FirstYearAttrition}.

\subsection{Intervention Practice Gaps: Documentation, Evaluation, and Experimental Frictions}
In many higher-education settings, intervention practices are deployed in an ad hoc manner, varying across instructors, courses, and academic terms, and are often insufficiently documented to support systematic evaluation \cite{Connolly2018TrialsEvidenceBased,Harackiewicz2018TargetedIntervention,Kizilcec2020ScalingInterventions}. The absence of structured intervention policies and consistent logging makes it difficult to determine when an intervention occurred, to whom it applied, and under which conditions \cite{Sonderlund2018LAInterventions,Page2020ClusteredObservationalEducation,Kitto2023CausalModelsLA}. As a result, routine observational data typically lack the granularity required to support credible causal evaluation of intervention effects in real-world academic operations \cite{Gallagher2020ChallengeBased,Sidiropoulos2022AgencySustainability}. Although randomized experiments (e.g., A/B testing) can, in principle, address these identification gaps, their implementation in live academic environments is frequently constrained by operational and practical considerations---such as coordination overhead, time-to-evidence, and uncertainty about where experimentation is most likely to yield actionable insights \cite{Chan2019AIMedEd,Connolly2018TrialsEvidenceBased,HahsVaughn2024MultisiteTrialChallenges}. In this setting, scenario-based structural comparison is not a substitute for causal inference, but a disciplined way to study how a fitted temporal model behaves under explicit intervention rules when intervention histories are not identifiable \cite{Igelstrom2022CausalInferenceObs,DoutreligneVaroquaux2025PredictiveDecisionCausal}.